\section{Sparse polynomial optimization}\label{sec:sp}

This section develops the sparse-bag route. The objective is a sum of
explicit polynomial factors of fixed degree on a bounded mixed box, and every
factor scope lies in a bag of a supplied tree decomposition with at most $p$
coordinates per bag. Every coordinate receives an independent linear
coefficient from one finite rational law chosen before sampling, that is,
the ambient model of \cref{def:model:perturbation} with marginal law
$U_{\sigma,M}$. No convexity, uniqueness, growth, nondegeneracy, or bound on
the integer dimension or the negative inertia is assumed.

The algorithm follows the common plan of the paper.
\begin{enumerate}[label=(\roman*)]
\item Curvature-corrected refinement of nested bag cells preserves every
global optimizer. The correction is one global rounding allowance
(\cref{sec:sp:cells}).
\item Independent linear noise bounds the expected number of retained bag
cells at every mesh by a quantity that does not depend on the mesh
(\cref{sec:sp:count}).
\item A certified local convexity test on a box that contains every global
optimizer finishes the sampled problem exactly (\cref{sec:sp:closure}).
\item A base-selected finite law makes failure of this test rare enough to
pay for an exact fallback on the same draw (\cref{sec:sp:law}).
\end{enumerate}
\Cref{thm:sp:main} states the result for fixed-degree polynomials on mixed
boxes. Quadratic and mixed-integer quadratic objectives are the case $d=2$;
their output is rational on every draw (\cref{cor:sp:qp}).
\Cref{sec:con} extends the method to feasible sets that are not boxes.

\subsection{Setting and main result}\label{sec:sp:setting}

\begin{definition}[Sparse mixed polynomial instance]\label{def:sp:instance}
A sparse mixed polynomial instance consists of the following rational data.
\begin{enumerate}[label=(\roman*)]
\item A mixed box $X=\prod_{i=1}^nX_i$ after the preprocessing of
\cref{sec:model:input}: $X_i=[\ell_i,u_i]$ for $i\in I_C$ and
$X_i=[\ell_i,u_i]\cap\Z$ with integers $\ell_i<u_i$ for $i\in I_Z$, where
$[n]=I_C\cup I_Z$ and every width $w_i=u_i-\ell_i$ is positive. Write
$\widehat X=\prod_i[\ell_i,u_i]$ for the continuous hull,
$W=\sum_iw_i$, $w_{\max}=\max_iw_i$, $n_c=\abs{I_C}$, $n_z=\abs{I_Z}$, and
$N_Z=\prod_{i\in I_Z}(w_i+1)$ for the number of integer assignments.
\item Factors $f_a$, $a\in\mathcal A$, each an explicit list of monomials
with rational coefficients, of total degree at most $d$, with variable scope
$S_a\subseteq[n]$. The objective is $F_0=\sum_af_a$.
\item A tree decomposition: a tree $\mathcal T$ whose nodes are bags
$\beta\subseteq[n]$, such that every coordinate and every scope $S_a$ lies
in some bag and, for every $i$, the bags containing $i$ form a subtree.
The bag size is $p=\max_\beta\abs\beta$.
\item A noise half-width $\sigma>0$ and a curvature bound $L>0$ with
\begin{equation}\label{eq:sp:curv}
 \partial_{ii}F_0(x)\le L\qquad\text{for all }x\in\widehat X\text{ and }i\in[n].
\end{equation}
\end{enumerate}
The base length $I$ is the binary length of these data. The degree $d\ge1$
is fixed.
\end{definition}

The curvature bound is needed on the whole continuous hull, including the
segments between adjacent integer labels, because the rounding and
comparison arguments move one coordinate along such segments. Inputs with
fixed coordinates reduce to this form by substitution, which does not
increase degrees, scopes or~\eqref{eq:sp:curv}; the sampled coefficient of a
substituted coordinate becomes an additive constant, which is carried
separately and added to reported values. If no coordinate remains, the
objective is a constant.

\paragraph{Status of the premises.}
The tree decomposition is verified in time polynomial in $I$. The bound
\eqref{eq:sp:curv} is promised: correctness of the pruning and of the patch
certificate relies on it. It can always be verified instead of promised.
With $R_0=\max\{1,\max_i\abs{\ell_i},\max_i\abs{u_i}\}$ and $F_0=\sum_\alpha
c_\alpha x^\alpha$,
\begin{equation}\label{eq:sp:Lmon}
 L_{\rm mon}=\max\Bigl\{1,\ \max_i\sum_\alpha\abs{c_\alpha}\,
 \alpha_i(\alpha_i-1)\,R_0^{\abs\alpha-2}\Bigr\}
\end{equation}
satisfies \eqref{eq:sp:curv} and has encoding length $\poly_d(I)$, because
$\abs{x^{\alpha-2e_i}}\le R_0^{\abs\alpha-2}$ on $\widehat X$ when
$\alpha_i\ge2$. Its numerical value can be much larger than a sharper
supplied bound, and it is the numerical value of $L$ that enters the work
bound. A sharper bound used in an independently checkable proof record must
be certified, with its verification cost added to the work. No general test
of polynomial inequalities is assumed. The same monomial method gives
rational numbers
\begin{equation}\label{eq:sp:kappa}
 \kappa_2\ge\max\Bigl\{1,\max_i\sum_k\sup_{\widehat X}\abs{\partial_{ik}F_0}\Bigr\},
 \qquad
 \kappa_3\ge\max\Bigl\{1,\max_i\sum_{k,l}\sup_{\widehat X}\abs{\partial_{ikl}F_0}\Bigr\}
\end{equation}
of encoding length $\poly_d(I)$, which the algorithm computes itself. Their
magnitudes affect only the stopping depth, logarithmically.

\paragraph{Noise and growth.}
The coefficients $\gamma_1,\dots,\gamma_n$ are independent with law
$U_{\sigma,M}$ of \cref{def:model:grid}, $F_\gamma=F_0+\gamma^{\mathsf T}x$,
and $f^*=\min_XF_\gamma$; the minimum exists because $X$ is compact. The
point-growth modulus $g_*(\gamma)$ of \cref{def:count:growth} is the largest
$g\ge0$ such that some minimizer $a$ satisfies
\begin{equation}\label{eq:sp:growth}
 F_\gamma(x)-F_\gamma(a)\ge g\norm{x-a}_2^2\qquad(x\in X).
\end{equation}
The set of admissible $g$ is closed, so the largest one exists. If
$g_*(\gamma)>0$ the minimizer is unique; if $F_\gamma$ has two distinct
minimizers, $g_*(\gamma)=0$.

\paragraph{Outputs.}
The algorithm returns an implicit patch output
(\cref{def:model:outputs}\ref{it:model:implicit}) on ordinary draws and the
algebraic output of the fallback of \cref{thm:count:fallback} on the
remaining draws. A patch output here is a quadruple $\Pi=(S,\bar x_S,P,g_0)$:
a set $S\supseteq I_Z$, rational values $\bar x_i\in X_i$ for $i\in S$
(integer labels, original continuous bounds, or values forced by a degenerate
hull interval), a rational box $P=\prod_{i\notin S}[a_i,b_i]$ with
$\ell_i\le a_i<b_i\le u_i$, which plays the role of the box $Q$ of
\cref{def:model:outputs}\ref{it:model:implicit}, and a rational $g_0>0$ for
which the run has verified
\begin{equation}\label{eq:sp:patch}
 \nabla^2_{\bar S\bar S}F_\gamma(\bar x_S,y)\succeq g_0I
 \qquad(y\in P),\qquad\bar S=[n]\setminus S.
\end{equation}
It denotes the unique minimizer $\hat y$ of $F_\gamma(\bar x_S,\cdot)$ on $P$,
and the run also proves that $\hat x=(\bar x_S,\hat y)$ is the unique global
minimizer of $F_\gamma$ on $X$. If $S=[n]$, the box $P$ is empty and
$\hat x=\bar x$. The patch output has length $\poly_d(I)$. The \emph{global
proof record} of the run consists of the candidate lists, incumbents and
cell bounds of every level, and the data of the closure tests; it proves
that every global minimizer lies in $\{\bar x_S\}\times P$.

\begin{theorem}[Sparse polynomial optimization on mixed boxes]\label{thm:sp:main}
Fix $d\ge1$. For sparse mixed polynomial instances
(\cref{def:sp:instance}) there is an algorithm and a power of two $M$,
computed from the base instance with $\log_2M\le\poly_d(I)$, such that for
the ambient perturbation model with marginal law $U_{\sigma,M}$ the
following hold.
\begin{enumerate}[label=(\roman*)]
\item\label{it:sp:exact} On every draw the algorithm returns an exact global
optimizer of $F_\gamma$ on $X$. On ordinary draws the output is a patch
output of length $\poly_d(I)$ together with its global proof record. On the
remaining draws, which have probability at most $1/(2B)$, it is the
algebraic output of \cref{thm:count:fallback} computed on the same draw,
where $B$ is a base-only factor with $\log B\le\poly_d(I)$.
\item\label{it:sp:work} The expected bit work, including the computation of
$M$ and the sampling, and the expected length and verification work of the
global proof record are at most
\begin{equation}\label{eq:sp:bound}
 C_0^p\Bigl[4+\frac{(1+n/2)Lw_{\max}}{2\sigma}\Bigr]^p\poly_d(I),
\end{equation}
where $C_0$ is an absolute constant.
\item\label{it:sp:eval} For every integer $q\ge0$, a $q$-bit evaluation of
the output returns a rational point $y\in X$ with
$\norm{y-\hat x}_2\le2^{-q}$, where $\hat x$ is the returned optimizer, and
rationals $a$ and $F_\gamma(y)$ with $a\le f^*\le F_\gamma(y)\le a+2^{-q}$.
It costs $\poly_d(I+q)$ bit operations for a patch output and
$B\poly_d(I+q)$ for a fallback output; its expected cost over the draw is
$\poly_d(I+q)$.
\end{enumerate}
The validity of $L$ is the only promised premise, and it is verified when
$L=L_{\rm mon}$. No growth, uniqueness or nondegeneracy premise is used; such
conditions enter only the analysis of the stopping depth.
\end{theorem}

The bound \eqref{eq:sp:bound} is polynomial for each fixed $p$ when
$Lw_{\max}/\sigma$ is polynomially bounded in $I$. It is a fixed-width
polynomial bound in the sense of \cref{sec:model:parameters}, not a
fixed-parameter bound in $p$: the bracket contains $n$, so the bound grows as
$n^p$ for fixed numerical scales. \Cref{sec:sp:width} explains where this
factor comes from, and \cref{thm:lim:width} shows that it is necessary for
the retention rule used here. The integer widths $w_i$ enter through their
magnitudes, so huge binary-encoded integer ranges are not covered at
polynomial cost; they enter the stopping depth, the sampling precision and
the fallback factor only logarithmically. The theorem concerns the sampled
objective. Its consequence for the base objective is the regret bound of
\cref{prop:model:regret} with $\omega_X(\gamma)\le\sigma W$. Finally, the
finite law matters: \cref{ex:model:tie} is an instance in this class on
which a tie occurs with probability $1/M$, so no patch certificate can
succeed on every draw and the fallback is necessary.

\subsection{Nested bag cells and sparse dynamic programming}\label{sec:sp:cells}

\paragraph{Grids and cells.}
Let $s$ be the least power of two with $s\ge\max\{1,w_{\max}\}$ and
$h_j=s2^{-j}$ for $j\ge0$. For $i\in I_C$, and for $i\in I_Z$ while
$h_j\ge1$, the level-$j$ nodes of coordinate $i$ are
\[
 G_{i,j}=\{\ell_i+kh_j:\ k\in\Z_{\ge0},\ \ell_i+kh_j<u_i\}\cup\{u_i\},
\]
and a \emph{coordinate cell} is $[a,a']\cap X_i$ for consecutive nodes
$a<a'$. For $i\in I_Z$ and $h_j<1$, the nodes are $G_{i,j}=X_i$ and the
coordinate cells are the singletons $\{a\}$, $a\in X_i$. While $h_j\ge1$,
$h_j$ is an integer, so integer nodes are integers. The \emph{endpoints} of
a cell are $a,a'$, or $a$ for a singleton. The level-$j$ grid is
$\mathbb G_j=\prod_iG_{i,j}\subseteq X$.

\begin{lemma}[Coordinate cells]\label{lem:sp:cells}
For every coordinate $i$ and level $j$:
\begin{enumerate}[label=(\alph*)]
\item the coordinate cells cover $X_i$, and a node contained in a cell is an
endpoint of it;
\item every node is an endpoint of at most two cells, and of exactly one if
the cells are singletons;
\item every cell has width at most $h_j$;
\item every level-$(j+1)$ cell lies in a level-$j$ cell, and every level-$j$
cell is the union of at most two level-$(j+1)$ cells.
\end{enumerate}
\end{lemma}

\begin{proof}
Consecutive nodes differ by at most $h_j$, and the first and last nodes are
$\ell_i$ and $u_i$; this gives (a) and (c) for interval cells, and (b) holds
because a node is the right endpoint of at most one cell and the left
endpoint of at most one cell. For (d), $h_{j+1}=h_j/2$ gives
$G_{i,j}\subseteq G_{i,j+1}$, and between consecutive level-$j$ nodes
$a<a'$ there is at most one level-$(j+1)$ node, namely $a+h_{j+1}$ if
$a+h_{j+1}<a'$. When $h_j=1>h_{j+1}$ for $i\in I_Z$, the level-$j$ cells
are $\{a,a+1\}$, the union of the singleton cells $\{a\}$ and $\{a+1\}$;
singleton cells are never refined further.
\end{proof}

A \emph{bag cell} of a bag $\beta$ at level $j$ is a product
$C=\prod_{i\in\beta}C_i$ of level-$j$ coordinate cells; its \emph{corners}
are the tuples of endpoints, and its \emph{children} are the level-$(j+1)$
bag cells contained in it. By \cref{lem:sp:cells}, a cell has at most
$2^{\abs\beta}$ corners and at most $2^{\abs\beta}$ children, its children
cover it, and a tuple of nodes is a corner of at most $2^{\abs\beta}$ cells.
The algorithm keeps for every bag a \emph{candidate list} $\mathcal W_j(\beta)$
of level-$j$ bag cells. It determines the \emph{physical domain}
\[
 \Omega_j=\{x\in X:\ \text{for every bag }\beta,\ x_\beta\in C
 \text{ for some }C\in\mathcal W_j(\beta)\},
\]
the \emph{row list} $\mathcal R_j(\beta)$, the set of corners of the cells in
$\mathcal W_j(\beta)$, and the \emph{allowed grid}
$A_j=\{y\in\mathbb G_j:\ y_\beta\in\mathcal R_j(\beta)\text{ for every }\beta\}$.
The physical domain is used only in the proofs; the algorithm stores cells
and rows.

\begin{lemma}[Allowed grid]\label{lem:sp:allowed}
$A_j=\mathbb G_j\cap\Omega_j$.
\end{lemma}

\begin{proof}
Corners of a cell lie in the cell, so $A_j\subseteq\Omega_j$. Conversely, if
$y\in\mathbb G_j\cap\Omega_j$ and $y_\beta\in C\in\mathcal W_j(\beta)$, each
$y_i$, $i\in\beta$, is a node contained in $C_i$, hence an endpoint of $C_i$
by \cref{lem:sp:cells}(a), so $y_\beta$ is a corner of $C$.
\end{proof}

\paragraph{Sparse dynamic programming.}
We may assume that every node of $\mathcal T$ has degree at most three:
a bag with children $\beta_1,\dots,\beta_k$, $k\ge3$, is replaced by a path
of $k-1$ copies of itself, the $l$th copy having child $\beta_l$ and the last
copy also $\beta_k$. Copies preserve the subtree property and the bag size
and at most double the number of nodes \citep{kloks1994-treewidth-computations-and-approximations}. Copies of one bag
have identical candidate lists throughout, because their cell bounds below
coincide. Assign every factor $f_a$ and every noise term $\gamma_ix_i$ to one
bag containing its scope, and let $\phi_\beta$ be the sum of the terms
assigned to $\beta$, so that $F_\gamma=\sum_\beta\phi_\beta(x_\beta)$. For a
row $v\in\mathcal R_j(\beta)$ let
\[
 m_\beta(v)=\min\{F_\gamma(y):\ y\in A_j,\ y_\beta=v\}
\]
be its \emph{min-marginal}, with $\min\emptyset=+\infty$, and let
$m_j=\min_{A_j}F_\gamma$.

\begin{lemma}[Sparse two-pass dynamic programming]\label{lem:sp:dp}
Let finite sets $G_i$, a tree decomposition of maximum degree three, finite
row sets $\mathcal R(\beta)\subseteq\prod_{i\in\beta}G_i$ and costs
$\phi_\beta:\mathcal R(\beta)\to\Q$ be given, and let
$A=\{y\in\prod_iG_i:y_\beta\in\mathcal R(\beta)\ \forall\beta\}$. Two passes
of messages compute every min-marginal
$m_\beta(v)=\min\{\sum_{\beta'}\phi_{\beta'}(y_{\beta'}):y\in A,\ y_\beta=v\}$,
the minimum over $A$, and a minimizer, using $O(\sum_\beta\abs{\mathcal R(\beta)})$
evaluations of the costs, additions, comparisons and dictionary operations
on separator tuples.
\end{lemma}

\begin{proof}
For adjacent bags $\beta,\beta'$ let $\Sigma=\beta\cap\beta'$, and let
$\mathcal T_{\beta\to\beta'}$ be the component of $\mathcal T$ minus the edge
$\beta\beta'$ that contains $\beta$. Define the message
$\mu_{\beta\to\beta'}$ on separator tuples $\varsigma$ by
\[
 \mu_{\beta\to\beta'}(\varsigma)=\min_{v\in\mathcal R(\beta),\,v_\Sigma=\varsigma}
 \Bigl[\phi_\beta(v)+\sum_{\beta''}\mu_{\beta''\to\beta}(v_{\beta''\cap\beta})\Bigr],
\]
the sum over the neighbors $\beta''\ne\beta'$ of $\beta$, and store only its
finite values in a dictionary keyed by $\varsigma$. By induction on the size
of $\mathcal T_{\beta\to\beta'}$, $\mu_{\beta\to\beta'}(\varsigma)$ is the
minimum of $\sum_{\beta''\in\mathcal T_{\beta\to\beta'}}\phi_{\beta''}$ over
the assignments to the coordinates of the bags of
$\mathcal T_{\beta\to\beta'}$ that equal $\varsigma$ on $\Sigma$ and whose
restriction to every bag of $\mathcal T_{\beta\to\beta'}$ is a row of that
bag. Indeed, the subtrees $\mathcal T_{\beta''\to\beta}$ partition
$\mathcal T_{\beta\to\beta'}$ minus $\beta$, and by the subtree property a
coordinate occurring in $\mathcal T_{\beta''\to\beta}$ and outside it lies in
$\beta''\cap\beta$; so once $v$ is fixed on $\beta$ the minimizations over
the different subtrees are independent and constrained only through their
separators. The same argument with all neighbors gives
$m_\beta(v)=\phi_\beta(v)+\sum_{\beta''}\mu_{\beta''\to\beta}(v_{\beta''\cap\beta})$,
the sum over all neighbors; every coordinate lies in some bag, so the
assignments are complete. Messages are computed toward a root and back,
recomputing each outgoing sum from at most two incoming messages instead of
subtracting, so that no difference of infinities occurs. Each row is
processed a bounded number of times, and stored argmin pointers give a
minimizer.
\end{proof}

This is nonserial dynamic programming on a tree decomposition
\citep{bertele1972-nonserial-dynamic-programming}, applied to sparse row lists. Dense tables
would enumerate products of coordinate grids; here each row is a corner of a
retained cell, and adjacent row lists are never multiplied. Every message
value is a sum of at most one assigned cost per bag, so its bit length is
polynomial in the bit length of the costs and the number of bags,
independently of the number of rows.

\paragraph{Rounding and pruning.}
The next lemma is the deterministic core of the method. It rounds all
coordinates at once and keeps the result inside every candidate list.

\begin{lemma}[Fixed-cell rounding]\label{lem:sp:round}
Let $x\in\Omega_j$, and let $C_0\in\mathcal W_j(\beta_0)$ be a cell with
$x_{\beta_0}\in C_0$. There is a random point $Y\in A_j$ such that
$Y_{\beta_0}$ is a corner of $C_0$, $Y_i=x_i$ whenever $x_i\in G_{i,j}$, and
\begin{equation}\label{eq:sp:round}
 \E F_\gamma(Y)\le F_\gamma(x)+E_j,\qquad E_j=\frac{nLh_j^2}{8}.
\end{equation}
\end{lemma}

\begin{proof}
For each $i$ with $x_i\in G_{i,j}$ put $Y_i=x_i$. Otherwise $x_i$ lies
strictly between consecutive nodes $a_i<a_i'$; this happens only for
$i\in I_C$ or for $i\in I_Z$ with $h_j\ge1$, so $a_i,a_i'\in X_i$. Let
$Y_i=a_i'$ with probability $(x_i-a_i)/(a_i'-a_i)$ and $Y_i=a_i$ otherwise,
independently over $i$. Then $\E Y_i=x_i$ and
$\operatorname{Var}Y_i=(x_i-a_i)(a_i'-x_i)\le h_j^2/4$.

Let $\beta$ be a bag and $C\in\mathcal W_j(\beta)$ a cell with $x_\beta\in
C$. For $i\in\beta$, if $x_i$ is a node then $Y_i=x_i\in C_i$ is an endpoint
of $C_i$ by \cref{lem:sp:cells}(a). Otherwise the only coordinate cell
containing $x_i$ is the one with endpoints $a_i,a_i'$, so $C_i$ is that cell
and $Y_i$ is an endpoint. Hence $Y_\beta$ is a corner of $C$. Since
$x\in\Omega_j$, every bag has such a cell, so $Y\in A_j$; taking
$\beta=\beta_0$, $C=C_0$ shows that $Y_{\beta_0}$ is a corner of $C_0$.

For the objective, put $Y^{(k)}=(Y_1,\dots,Y_k,x_{k+1},\dots,x_n)$. Given
$Y_1,\dots,Y_{k-1}$, the function
$\varphi(\theta)=F_\gamma(Y_1,\dots,Y_{k-1},\theta,x_{k+1},\dots,x_n)$ on
$[\ell_k,u_k]$ satisfies $\varphi''\le L$ by \eqref{eq:sp:curv}, because
these points lie in $\widehat X$ and the linear noise does not change second
derivatives. Taylor's theorem gives
$\varphi(Y_k)\le\varphi(x_k)+\varphi'(x_k)(Y_k-x_k)+\frac L2(Y_k-x_k)^2$.
As $Y_k$ is independent of $Y_1,\dots,Y_{k-1}$ with mean $x_k$,
\[
 \E[F_\gamma(Y^{(k)})\mid Y_1,\dots,Y_{k-1}]\le F_\gamma(Y^{(k-1)})
 +\frac L2\operatorname{Var}Y_k .
\]
Taking expectations and summing over $k$ gives \eqref{eq:sp:round}.
\end{proof}

For a quadratic, \eqref{eq:sp:round} also follows from cancellation of the
off-diagonal terms in expectation. The sequential argument is needed for
higher degree, for instance for the monomial $x_1^2x_2^2$, whose expected
rounding error is not a sum of diagonal terms.

For a cell $C\in\mathcal W_j(\beta)$ define its \emph{cell bound}
\[
 q_C=\min\{m_\beta(v):\ v\text{ a corner of }C\}.
\]
The \emph{incumbent} $U_j$ is the least value of $F_\gamma$ at the allowed
grid points produced by the dynamic programs of levels $0,\dots,j$; it is the
value of a point of $X$, and $U_j\le m_j$. The candidate list is updated by
the \emph{retention rule}
\begin{equation}\label{eq:sp:retain}
 \mathcal W'_j(\beta)=\{C\in\mathcal W_j(\beta):\ q_C-E_j\le U_j\},\qquad
 \mathcal W_{j+1}(\beta)=\{\text{children of cells in }\mathcal W'_j(\beta)\},
\end{equation}
starting from $\mathcal W_0(\beta)=\{X_\beta\}$. Ties are retained. Let
$\Omega'_j$ be the physical domain of the retained lists; it equals
$\Omega_{j+1}$, because children cover their parents.

\begin{proposition}[Sound pruning and witnesses]\label{prop:sp:prune}
On every draw and at every level $j\ge0$:
\begin{enumerate}[label=(\alph*)]
\item every global minimizer $x^*$ of $F_\gamma$ on $X$ lies in $\Omega_j$,
and $x^*_\beta$ lies in a retained cell of every bag;
\item $f^*\le U_j\le m_j\le f^*+E_j$;
\item every $x\in\Omega_j\setminus\Omega'_j$ satisfies $F_\gamma(x)>U_j$;
\item every retained cell $C\in\mathcal W'_j(\beta)$ has a corner $v$ and a
point $y\in A_j$ with $y_\beta=v$ and $F_\gamma(y)=q_C\le f^*+2E_j$.
\end{enumerate}
\end{proposition}

\begin{proof}
We first note the cell lower bound. Let $C\in\mathcal W_j(\beta)$ and let
$x\in\Omega_j$ have $x_\beta\in C$. \Cref{lem:sp:round} with $(\beta,C)$ gives
$Y\in A_j$ with $Y_\beta$ a corner of $C$, so
$F_\gamma(Y)\ge m_\beta(Y_\beta)\ge q_C$ on every outcome, and
\begin{equation}\label{eq:sp:cellbound}
 q_C-E_j\le F_\gamma(x).
\end{equation}

We prove (a) by induction; $\Omega_0=X$. Suppose every minimizer lies in
$\Omega_j$, and let $x^*$ be one. For each bag $\beta$ choose
$C\in\mathcal W_j(\beta)$ containing $x^*_\beta$. By \eqref{eq:sp:cellbound},
$q_C-E_j\le f^*\le U_j$, so $C$ is retained, and its children cover it; hence
$x^*\in\Omega_{j+1}$. For (b), $U_j\ge f^*$ because $U_j$ is a value at a
point of $X$, and rounding a minimizer by \cref{lem:sp:round} gives an
allowed point of expected value at most $f^*+E_j$, so $m_j\le f^*+E_j$.
For (c), $x\in\Omega_j\setminus\Omega'_j$ lies in some cell
$C\in\mathcal W_j(\beta)$ but in no retained cell of $\beta$, so $C$ was
deleted, and \eqref{eq:sp:cellbound} gives $F_\gamma(x)\ge q_C-E_j>U_j$.
For (d), $q_C\le U_j+E_j<\infty$, so $q_C=m_\beta(v)$ for some corner $v$,
and the minimum defining $m_\beta(v)$ is attained at some $y\in A_j$. By (b),
$q_C\le U_j+E_j\le f^*+2E_j$.
\end{proof}

Part~(a) holds for tied and positive-dimensional optimal sets and for every
draw. Part~(d) gives one globally consistent feasible point for each retained
cell; the counting argument below needs exactly this, and a collection of
separate per-bag witnesses would not suffice. Part~(c) is the global
certificate: a point removed at level $j$ has value above $U_j\ge f^*$.
An incumbent found at an earlier level may lie outside the current physical
domain; it is still feasible and remains a valid upper bound.

\subsection{Expected size of the retained lists}\label{sec:sp:count}

The count conditions on the noise outside one bag. For a bag $\beta$, write
$O=[n]\setminus\beta$ and, for $v\in\widehat X_\beta=\prod_{i\in\beta}[\ell_i,u_i]$,
\begin{equation}\label{eq:sp:V}
 V_\beta(v)=\min_{z\in X_O}\bigl[F_0(v,z)+\gamma_O^{\mathsf T}z\bigr].
\end{equation}
The minimum is over the original outside domain $X_O=\prod_{i\in O}X_i$, not
over a set that depends on the run. Hence $V_\beta$ depends on $\gamma_O$ and
not on $\gamma_\beta$, and $f^*=\min_{v\in X_\beta}[V_\beta(v)+\gamma_\beta^{\mathsf T}v]$.

\begin{lemma}[Comparison intervals]\label{lem:sp:compare}
Fix $\gamma_O$, a bag tuple $v\in X_\beta$ and $\eta\ge0$, and suppose
$V_\beta(v)+\gamma_\beta^{\mathsf T}v\le f^*+\eta$. If $i\in\beta$ and
$a>0$ satisfy $v\pm ae_i\in X_\beta$, then $\gamma_i$ lies in a closed
interval that depends only on $\gamma_O$, $v$, $i$, $a$ and $\eta$, and whose
length is at most $La+2\eta/a$.
\end{lemma}

\begin{proof}
Let $z^*$ attain the minimum in \eqref{eq:sp:V} at $v$. For
$\theta\in[-a,a]$ put $\psi(\theta)=F_0(v+\theta e_i,z^*)$; the segment lies
in $\widehat X$, so $\psi''\le L$ and $\psi(a)+\psi(-a)-2\psi(0)\le La^2$.
Since $V_\beta(v\pm ae_i)\le\psi(\pm a)+\gamma_O^{\mathsf T}z^*$ and
$V_\beta(v)=\psi(0)+\gamma_O^{\mathsf T}z^*$,
\[
 V_\beta(v+ae_i)+V_\beta(v-ae_i)-2V_\beta(v)\le La^2 .
\]
The tuples $v\pm ae_i$ are feasible, so
$f^*\le V_\beta(v\pm ae_i)+\gamma_\beta^{\mathsf T}(v\pm ae_i)$. Subtracting
the hypothesis gives
\[
 \frac{V_\beta(v)-V_\beta(v+ae_i)-\eta}{a}\le\gamma_i\le
 \frac{V_\beta(v-ae_i)-V_\beta(v)+\eta}{a}.
\]
The endpoints do not involve $\gamma_\beta$, and the difference of the
endpoints is at most $La+2\eta/a$ by the previous display.
\end{proof}

For a bag $\beta$ and level $j$ let
\[
 N_j(\beta)=\bigl\{v\in\mathbb G_{j,\beta}:\
 V_\beta(v)+\gamma_\beta^{\mathsf T}v\le f^*+2E_j\bigr\},\qquad
 \mathbb G_{j,\beta}=\prod_{i\in\beta}G_{i,j},
\]
a subset of the complete deterministic bag grid. The algorithm never
enumerates it; it serves only as a dominating count.

\begin{proposition}[Expected retained cells]\label{prop:sp:count}
Let $J\ge0$, let $M\ge2^J$, and let $0\le j\le J$. On every draw,
$\abs{\mathcal W'_j(\beta)}\le2^{\abs\beta}\abs{N_j(\beta)}$. Moreover
\begin{equation}\label{eq:sp:Theta}
 \E\abs{N_j(\beta)}\le\Theta_\beta:=\prod_{i\in\beta}
 \Bigl[4+\frac{(1+n/2)Lw_i}{2\sigma}\Bigr].
\end{equation}
\end{proposition}

\begin{proof}
Let $C\in\mathcal W'_j(\beta)$ and let $v$, $y$ be as in
\cref{prop:sp:prune}(d). Since $y=(v,y_O)$ with $y_O\in X_O$,
$V_\beta(v)+\gamma_\beta^{\mathsf T}v\le F_\gamma(y)\le f^*+2E_j$, so
$v\in N_j(\beta)$. A node tuple is a corner of at most $2^{\abs\beta}$
level-$j$ cells, which gives the first claim.

For the expectation, let $a_i=h_j$ if $i\in I_C$ or $h_j\ge1$, and $a_i=1$
otherwise; thus $a_i\ge h_j$, and consecutive nodes of coordinate $i$ are
$a_i$ apart, except possibly the last two. Call a node $v_i\in G_{i,j}$
\emph{regular} if $v_i\pm a_i\in X_i$. Regular nodes are $\ell_i+ka_i$ with
$k\ge1$ and $\ell_i+(k+1)a_i\le u_i$, so there are at most $w_i/a_i$ of
them, and at most three nodes are not regular: $\ell_i$, the last node
below $u_i$, and $u_i$.

Fix $v\in\mathbb G_{j,\beta}$ and condition on $\gamma_O$. If
$v\in N_j(\beta)$, \cref{lem:sp:compare} with $\eta=2E_j$ confines every
$\gamma_i$, $i\in\beta$ with $v_i$ regular, to an interval determined by
$\gamma_O$ and $v$, of length at most $\lambda_i=La_i+4E_j/a_i$. The
coefficients $\gamma_i$, $i\in\beta$, are independent of each other and of
$\gamma_O$, so by \eqref{eq:model:interval}
\[
 \Prob\{v\in N_j(\beta)\mid\gamma_O\}\le\prod_{i\in\beta}c_i(v_i),
 \qquad c_i(v_i)=\begin{cases}\lambda_i/(2\sigma)+1/M&v_i\text{ regular},\\
 1&\text{otherwise}.\end{cases}
\]
The right side is a product of functions of single coordinates, so summing
over the grid factorizes:
\[
 \E\abs{N_j(\beta)}\le\prod_{i\in\beta}\sum_{v_i\in G_{i,j}}c_i(v_i)
 \le\prod_{i\in\beta}\Bigl[3+\frac{w_i}{a_i}\Bigl(\frac{\lambda_i}{2\sigma}
 +\frac1M\Bigr)\Bigr].
\]
Here $\frac{w_i\lambda_i}{2\sigma a_i}=\frac{Lw_i}{2\sigma}\bigl(1+\frac{nh_j^2}{2a_i^2}\bigr)
\le\frac{(1+n/2)Lw_i}{2\sigma}$ because $a_i\ge h_j$, and
$\frac{w_i}{a_iM}\le\frac{s}{h_jM}=\frac{2^j}M\le1$ because $j\le J$ and
$M\ge2^J$. This proves \eqref{eq:sp:Theta}.
\end{proof}

The expectation is taken over the complete deterministic bag grid, and every
retained cell is charged to a near-optimal node of that grid on every draw.
The argument therefore does not assume that the adaptively retained cells
are independent of the noise. It is the bag-conditional analogue of
\cref{thm:count:local} and \cref{thm:count:cells}(iv). The proof above is
self-contained because the grids here are clipped at the upper bounds and the
count is conditioned on the noise outside the bag; the clipped nodes are
treated as exceptional and are never used in comparisons. The requirement
$M\ge2^J$ is what makes the count independent of the level: for a fixed
coarse grid, the atomic term
$w_i/(a_iM)$ would grow without bound as $h_j\to0$.

\subsection{Why the bound is fixed-width but not fixed-parameter}\label{sec:sp:width}

The tree decomposition controls the shape of the dynamic program: tables
are indexed by bag rows, messages by separator tuples, and the number of
rows of a bag is charged to $\Theta_\beta$, a product over the coordinates
of that bag only. Treewidth therefore governs the exponent in
\eqref{eq:sp:bound}. The base of the exponential, however, contains the
total dimension $n$. It enters through the rounding allowance: the cell bound
\eqref{eq:sp:cellbound} must hold for every feasible point of a cell, and
its proof rounds all $n$ coordinates of a global point simultaneously, so
the allowance $E_j=nLh_j^2/8$ is a sum of variances over all coordinates.
The witnesses of \cref{prop:sp:prune}(d) are near-optimal only up to
$2E_j$, and \cref{lem:sp:compare} turns this tolerance into noise intervals
of length $(1+n/2)La_i$ for every coordinate of the bag. The resulting bound
on the expected number of retained rows of a bag is of order
$(nLw_{\max}/\sigma)^{\abs\beta}$ when this ratio is large, which is polynomial
in $n$ only for fixed $p$.

Two facts show that this is a property of the method and not of the
analysis. First, replacing $E_j$ by an allowance for the bag alone,
$\abs\beta Lh_j^2/8$, is unsound: \cref{prop:lim:local} gives a star
quadratic of width two on which the local allowance deletes the cell that
contains every global optimizer. The error that the global allowance absorbs
comes from minimizing over the grid outside the bag; it varies with the
bag's own coordinates and does not cancel when min-marginals are compared.
Second, for the retention rule \eqref{eq:sp:retain} and closure on the whole
hull, the factor $n^{\Omega(p)}$ actually occurs: \cref{thm:lim:width} gives
connected instances of treewidth $p-1$, with $L=2$, unit widths and
$\sigma=1/100$, on which a bag retains more than $(5n/(6p))^{p/2}$ cells on
every draw before the stopping depth is reached. A dimension-free count would
require exact or certified conditional values for the variables outside a
bag; \cref{sec:rec} develops such a conditional oracle for a small core with
tractable recourse, which is a different structural assumption.

\subsection{Certified closure on a strongly convex patch}\label{sec:sp:closure}

After the retention step of level $j$, the algorithm tries to finish. Let
$g_0>0$ be a rational number chosen from the base data (\cref{def:sp:schedule})
and let $\kappa_2,\kappa_3$ be as in \eqref{eq:sp:kappa}.

\begin{definition}[Closure test]\label{def:sp:closure}
At level $j$:
\begin{enumerate}[label=(\arabic*)]
\item For every coordinate $i$ let $Q_i$ be the intersection, over the bags
$\beta\ni i$, of the smallest closed interval containing the $i$th
projections of the cells of $\mathcal W'_j(\beta)$, and let
$Q=\prod_iQ_i$.
\item Record $\bar x_i$ equal to the point of $Q_i$ for every $i\in I_Z$ with
$Q_i$ a single point.
\item Let $c$ be the midpoint of $Q$ and $r$ the largest half-width of the
intervals $Q_i$. For $i\in I_C$, record $\bar x_i=\ell_i$ if
$\partial_iF_\gamma(c)-\kappa_2r>0$, and $\bar x_i=u_i$ if
$\partial_iF_\gamma(c)+\kappa_2r<0$.
\item If some integer coordinate is unrecorded, the test fails. Otherwise
record $\bar x_i$ equal to the point of $Q_i$ for every unrecorded $i\in I_C$
with $Q_i$ a single point. Let $S$ be the recorded set and $\bar S=[n]\setminus S$.
If $\bar S=\emptyset$, return the point $\bar x$.
\item Otherwise let $P=\prod_{i\in\bar S}Q_i$, with midpoint $c^P$ and
largest half-width $r^P$. If the rational symmetric matrix
\begin{equation}\label{eq:sp:matrix}
 \nabla^2_{\bar S\bar S}F_\gamma(\bar x_S,c^P)-(\kappa_3r^P+g_0)I
\end{equation}
is positive definite, return the patch output $(S,\bar x_S,P,g_0)$;
otherwise the test fails.
\end{enumerate}
\end{definition}

Positive definiteness of a rational symmetric matrix is decided exactly by
Gaussian elimination: all pivots must be positive. All quantities in the test
have encoding length polynomial in $I$, $j$ and $\log M$.

\begin{proposition}[Soundness of closure]\label{prop:sp:closure-sound}
On every draw, every equality recorded by the closure test holds at every
global minimizer of $F_\gamma$ on $X$. If the test returns a point $\bar x$,
then $\bar x$ is the unique global minimizer. If it returns
$\Pi=(S,\bar x_S,P,g_0)$, then \eqref{eq:sp:patch} holds on $P$, the function
$F_\gamma(\bar x_S,\cdot)$ has a unique minimizer $\hat y$ on $P$, and
$(\bar x_S,\hat y)$ is the unique global minimizer of $F_\gamma$ on $X$.
\end{proposition}

\begin{proof}
By \cref{prop:sp:prune}(a), every minimizer $x^*$ has $x^*_\beta$ in a
retained cell of every bag, so $x^*\in Q$. A singleton $Q_i$ therefore fixes
$x_i^*$; for $i\in I_Z$ its point is an endpoint of an integer cell, hence an
integer.

For the gradient test, let $x\in Q\subseteq\widehat X$. The mean value
theorem along the segment from $c$ to $x$, which lies in $Q$, gives
$\abs{\partial_iF_\gamma(x)-\partial_iF_\gamma(c)}\le\sum_k\sup_{\widehat X}
\abs{\partial_{ik}F_0}\abs{x_k-c_k}\le\kappa_2r$; the linear noise does not
change second derivatives. If $\partial_iF_\gamma(c)-\kappa_2r>0$, then
$\partial_iF_\gamma(x^*)>0$. Were $x^*_i>\ell_i$, the point
$x^*-\theta e_i$ would be feasible for small $\theta>0$, because $i$ is
continuous, and would have a smaller value. Hence $x^*_i=\ell_i$. The upper
bound is symmetric.

All recorded values lie in $X_i$, and $P\subseteq\prod_{i\in\bar S}[\ell_i,u_i]$
with $\bar S\subseteq I_C$, so $T=\{\bar x_S\}\times P\subseteq X$, and every
minimizer lies in $T$. If $\bar S=\emptyset$, a minimizer exists by
compactness and equals $\bar x$. Otherwise, for $y\in P$ every entry of
$\nabla^2_{\bar S\bar S}F_\gamma(\bar x_S,y)-\nabla^2_{\bar S\bar S}F_\gamma(\bar x_S,c^P)$
in row $k$ is at most $r^P\sum_m\sup_{\widehat X}\abs{\partial_{klm}F_0}$ in
absolute value, by the mean value theorem; the spectral norm of a symmetric
matrix is at most its largest absolute row sum, so the difference has norm
at most $\kappa_3r^P$. Positive definiteness of \eqref{eq:sp:matrix} gives
\eqref{eq:sp:patch}. Then $F_\gamma(\bar x_S,\cdot)$ is strictly convex on
$P$ and has a unique minimizer $\hat y$. Since $T\subseteq X$ contains every
global minimizer and at least one, $\min_TF_\gamma=f^*$ and the global
minimizers are exactly the minimizers on $T$, that is, $(\bar x_S,\hat y)$.
\end{proof}

The test fixes only original bounds and integer values that every
minimizer shares; it never fixes a coordinate to an artificial endpoint of
the hull, and it applies the gradient test only to continuous coordinates.
An interior integer optimum can have a nonzero derivative, and a positive
definite continuous block does not justify closure while integer choices
remain; both rules are therefore essential.

The test succeeds early when the sampled instance is well conditioned. The
conditions in the next proposition are not tested by the algorithm; they
are used only to bound the stopping depth.

\begin{proposition}[When closure succeeds]\label{prop:sp:closure-stop}
Suppose $g_*(\gamma)\ge g_0>0$, let $a$ be the unique minimizer, and suppose
that $\abs{\partial_iF_\gamma(a)}>\tau$ for every $i\in I_C$ with
$a_i\in\{\ell_i,u_i\}$. Let $A=2+nL/g_0$. If
\begin{equation}\label{eq:sp:cap}
 h_j\le\min\Bigl\{\frac1{4A},\ \frac{\tau}{4\kappa_2A},\
 \frac{g_0}{4\kappa_3A}\Bigr\},
\end{equation}
then the closure test of level $j$ returns a point or a patch output.
\end{proposition}

\begin{proof}
Let $C\in\mathcal W'_j(\beta)$, and let $y$ be its witness from
\cref{prop:sp:prune}(d). Growth gives $g_0\norm{y-a}^2\le F_\gamma(y)-f^*\le
2E_j=nLh_j^2/4$, so $\norm{y-a}_2\le\frac{h_j}2\sqrt{nL/g_0}$. Every point of
$C_i$ is within $h_j$ of $y_i$, so within
$h_j(1+\frac12\sqrt{nL/g_0})\le Ah_j$ of $a_i$. Since $a_i\in Q_i$, every
$Q_i$ lies in $[a_i-Ah_j,a_i+Ah_j]$, and $r\le Ah_j$ and
$\norm{c-a}_\infty\le r$.

By \eqref{eq:sp:cap}, $h_j\le1/(4A)<1$, so integer cells are singletons.
For $i\in I_Z$, the integer cell of $C$ is $\{y_i\}$ and
$\abs{y_i-a_i}\le Ah_j\le1/4$, so $y_i=a_i$. Every retained cell of every bag
containing $i$ therefore has $i$th projection $\{a_i\}$, $Q_i=\{a_i\}$, and
step~(2) records every integer coordinate.

Let $i\in I_C$ with $a_i=\ell_i$; then $\partial_iF_\gamma(a)>\tau$ by
hypothesis and first-order optimality. As in the proof of
\cref{prop:sp:closure-sound},
$\partial_iF_\gamma(c)-\kappa_2r\ge\partial_iF_\gamma(a)-2\kappa_2Ah_j>\tau-\tau/2>0$,
so step~(3) records $\bar x_i=\ell_i$; the case $a_i=u_i$ is symmetric. By
soundness no coordinate with $\ell_i<a_i<u_i$ is recorded at a bound. Hence
$\bar S$ consists of continuous coordinates that are interior at $a$. If
$\bar S=\emptyset$ the test returns a point. Otherwise, for every unit vector
$e$ supported on $\bar S$, the points $a\pm\theta e$ are feasible for small
$\theta>0$, and adding the two growth inequalities and letting
$\theta\to0$ gives $e^{\mathsf T}\nabla^2F_\gamma(a)e\ge2g_0$; thus
$\nabla^2_{\bar S\bar S}F_\gamma(a)\succeq2g_0I$. The point $(\bar x_S,c^P)$
differs from $a$ only in $\bar S$, by at most $r^P\le Ah_j$, so the matrix
\eqref{eq:sp:matrix} is at least
$(2g_0-2\kappa_3Ah_j-g_0)I\succeq\frac{g_0}2I$.
\end{proof}

The two-sided Taylor argument uses only the coordinates that are interior
at the optimizer. A boundary optimizer whose full Hessian is indefinite is
therefore covered: the active coordinates are removed before the matrix
test. No lower bound on the distance of the optimizer from inactive bounds
is needed.

\subsection{One finite law and the same-draw fallback}\label{sec:sp:law}

Two events make the closure test late: small point growth, and an active
continuous coordinate with a small derivative. Both have small probability
under the grid law.

\begin{lemma}[Finite-law tails on mixed boxes]\label{lem:sp:tails}
There is a base-computable integer $C_{\rm tail}\ge1$ with
$\log C_{\rm tail}\le\poly_d(I)$ such that, for every $M$ and every
$\varepsilon>0$,
\begin{equation}\label{eq:sp:growthtail}
 \Prob\{g_*<\varepsilon\}\le\frac{\varepsilon W}{\sigma}+\frac{2nC_{\rm tail}}M .
\end{equation}
Let $D=\max\{1,d-1\}$ and $K=\max\{1,n_c3^{n_c}N_ZD^{n_c}\}$. For every
$\tau>0$, the probability that $g_*>0$ and the unique minimizer $a$ has a
coordinate $i\in I_C$ with $a_i\in\{\ell_i,u_i\}$ and
$\abs{\partial_iF_\gamma(a)}\le\tau$ is at most $K(\tau/\sigma+1/M)$.
\end{lemma}

\begin{proof}
This is \cref{lem:count:finite-tails}(b) and~(d) with $m=0$, $S=W$ and
$R_Z=N_Z$; the first part combines \cref{thm:count:growth-tail} for the
continuous uniform law with the transfer of
\cref{lem:count:finite-tails}(a). The mixed box is described by $2n_c$
linear inequalities and by disjunctions over the $N_Z$ integer labels, so
the number $s$ of atoms satisfies $\log(s+1)\le\poly(I)$, and the constant
$C_{\rm tail}$ of \cref{lem:count:finite-tails} has logarithm $\poly_d(I)$
even though the number of labels can be exponential.
\end{proof}

The second bound counts nonsingular stationary points of faces and does not
assume that stationary sets are finite. It is an unconditional estimate for
the intersection with $\{g_*>0\}$, not a conditional probability given good
growth.

\begin{definition}[Base schedule]\label{def:sp:schedule}
Before sampling, compute from the base instance:
\begin{enumerate}[label=(\arabic*)]
\item $\kappa_2,\kappa_3$ by \eqref{eq:sp:kappa}, $C_{\rm tail}$ and $K$ as
in \cref{lem:sp:tails}, and the factor $B\ge2$ of \cref{thm:count:fallback}
for the format of the instance (the mixed box $X$, the degree $d$, and the
monomial support of $F_0$ with $n$ added linear terms);
\item $\rho=1/(4B)$, $g_0=\rho\sigma/(2W)$, $\tau=\rho\sigma/(2K)$ and
$A=2+nL/g_0$;
\item the stopping depth $J$, the least $j\ge0$ satisfying \eqref{eq:sp:cap};
\item the grid size $M$, the least power of two with
$M\ge\max\{2,\,2^J,\,4nC_{\rm tail}/\rho,\,2K/\rho\}$.
\end{enumerate}
\end{definition}

Every quantity in the schedule has encoding length $\poly_d(I)$: the
logarithms of $B$, $C_{\rm tail}$, $K$, $\kappa_2$, $\kappa_3$, $1/\rho$,
$1/g_0$, $1/\tau$ and $A$ are polynomial in $I$, so $J\le\poly_d(I)$ and
$\log_2M\le\poly_d(I)$. The order matters. The factor $B$ is fixed first;
the thresholds and $J$ depend on $B$ but not on $M$; and $M$ is chosen last.
The sampled coefficients have bit length $b=\log_2M+O(\operatorname{size}\sigma)$,
which enters only the polynomial factor of the fallback.

\begin{lemma}[Rare late closure]\label{lem:sp:rare}
Under the schedule of \cref{def:sp:schedule}, the closure test succeeds at
some level $j\le J$ except on an event of probability at most $1/(2B)$.
\end{lemma}

\begin{proof}
By \cref{lem:sp:tails},
$\Prob\{g_*<g_0\}\le\rho/2+2nC_{\rm tail}/M\le\rho$, and the probability
that $g_*\ge g_0$ while an active continuous derivative has magnitude at
most $\tau$ is at most $K\tau/\sigma+K/M\le\rho/2+\rho/2$. Outside these two
events \cref{prop:sp:closure-stop} applies at level $J$. The union has
probability at most $2\rho=1/(2B)$.
\end{proof}

\begin{definition}[Sparse cell algorithm]\label{def:sp:algorithm}
Given a sparse mixed polynomial instance:
\begin{enumerate}[label=(\arabic*)]
\item Compute the schedule of \cref{def:sp:schedule}, convert the tree
decomposition to maximum degree three, assign factors and noise terms to
bags, and draw $\gamma$ once from the grid law with this $M$.
\item Set $\mathcal W_0(\beta)=\{X_\beta\}$ for every bag. For
$j=0,1,\ldots,J$: form the row lists, run the dynamic program of
\cref{lem:sp:dp} for all min-marginals and a minimizer of $F_\gamma$ on
$A_j$, update the incumbent, apply the retention rule \eqref{eq:sp:retain},
and run the closure test of \cref{def:sp:closure}. If the test returns, stop
with its output. Otherwise form $\mathcal W_{j+1}$ by generating children and
removing duplicates.
\item If no closure test has returned by level $J$, run the exact fallback of
\cref{thm:count:fallback} on $F_\gamma$ over $X$ and return its output.
\end{enumerate}
\end{definition}

\begin{proof}[Proof of \cref{thm:sp:main}]
\emph{Exactness.} A closure output is correct on every draw by
\cref{prop:sp:closure-sound}, and the fallback is correct on every draw by
\cref{thm:count:fallback}. The fallback runs only if no closure test returns,
which by \cref{lem:sp:rare} has probability at most $1/(2B)$. No draw is
discarded or repeated.

\emph{Work of one level.} Let $j\ge1$. Generating $\mathcal W_j(\beta)$ costs
$O(2^{\abs\beta}\abs{\mathcal W'_{j-1}(\beta)})$ cell operations and a sort,
the row list has at most $2^{\abs\beta}\abs{\mathcal W_j(\beta)}$ entries, the
dynamic program costs $O(\sum_\beta\abs{\mathcal R_j(\beta)})$ operations by
\cref{lem:sp:dp}, and the retention rule and the hulls cost
$O(2^{\abs\beta}\abs{\mathcal W_j(\beta)})$ further operations; level~$0$ has
one cell per bag. Every grid coordinate at level $j\le J$ has bit length
$O(I+j+\log s)$, every sampled coefficient $O(I+\log M)$, every factor value
and message $\poly_d(I)$, and a sort or dictionary operation on a list of
$R$ rows costs $O(\log R)$ comparisons with $\log R\le\log\abs{\mathbb G_j}
\le\poly_d(I)$. The closure test costs $\poly_d(I)$ plus the hull
computation. Hence the work of level $j$ is at most
$8^p\sum_\beta(1+\abs{\mathcal W'_{j-1}(\beta)})\poly_d(I)$.

\emph{Expectation.} By \cref{prop:sp:count}, with $M\ge2^J$,
$\E\abs{\mathcal W'_{j-1}(\beta)}\le2^p\Theta_\beta$ for $1\le j\le J$, and
$\Theta_\beta\le[4+(1+n/2)Lw_{\max}/(2\sigma)]^p$. The work is a sum of
nonnegative terms, each bounded by a constant times a retained-list size,
so its expectation is bounded by linearity; no product of random table
sizes occurs. There are $J+1\le\poly_d(I)$ levels and polynomially many
bags. The fallback costs $B\poly_d(I+b)$ and runs with probability at most
$1/(2B)$, so it contributes $\poly_d(I)$ to the expectation
(\cref{lem:count:rare-fallback}). Computing the schedule and drawing
$n\log_2M$ random bits cost $\poly_d(I)$. This proves \eqref{eq:sp:bound}
with $C_0=16$. The global proof record consists of quantities computed
during the run, so its length is bounded by the work; it is verified by
recomputing the dynamic programs of the levels, which costs the same.

\emph{Evaluation.} Part~\ref{it:sp:eval} is \cref{prop:sp:eval} below.
\end{proof}

\subsection{Exact implicit output and arbitrary-precision evaluation}\label{sec:sp:eval}

A patch output determines one point exactly, but the coordinates of that
point are usually irrational algebraic numbers whose minimal polynomials can
have large degree. The output contract does not include them. What it
includes is evaluation to any requested precision in time polynomial in the
precision.

\begin{proposition}[Evaluation]\label{prop:sp:eval}
Let $q\ge0$.
\begin{enumerate}[label=(\alph*)]
\item Given a patch output $\Pi=(S,\bar x_S,P,g_0)$ of \cref{thm:sp:main}, a
rational $y\in P$ and a rational $a$ with
$a\le f^*\le F_\gamma(\bar x_S,y)\le a+2^{-q}$ and
$\norm{y-\hat y}_2\le2^{-q}$ are computed in $\poly_d(I+q)$ bit operations.
\item Given a fallback output, a rational $y\in X$ with the same integer
labels as the returned optimizer $x^\circ$, and a rational $a$ with
$a\le f^*\le F_\gamma(y)\le a+2^{-q}$ and $\norm{y-x^\circ}_2\le2^{-q}$, are
computed in $B\poly_d(I+q)$ bit operations.
\end{enumerate}
In particular the expected cost of a $q$-bit evaluation is $\poly_d(I+q)$.
\end{proposition}

The proof is in \cref{app:sp:eval}. Part~(a) applies weak convex
optimization in the bit model \citep[Definition~2.1.10 and
Corollary~4.2.7]{grotschel1988-geometric-algorithms-and-combinatorial-optimization} to the capped epigraph of
$F_\gamma(\bar x_S,\cdot)$ over $P$, with an explicit repair of the
approximately feasible output; the modulus $g_0$ converts value accuracy
into distance, including at boundary minimizers, and enters the cost only
through its encoding length. The cost does not depend on the numerical ratio
$L/g_0$, as a grid-based evaluation would. Evaluation does not decide exact
comparisons of $f^*$ with a rational threshold, and it does not identify the
active set of $\hat y$. Without a verified modulus, value accuracy alone
would not localize the optimizer: a function that is flat along a direction
has points of equal value far from each other.

\subsection{Quadratic and mixed-integer quadratic objectives}\label{sec:sp:qp}

For $d=2$ the patch and the fallback can both be solved in rational
arithmetic, so the output is rational on every draw.

\begin{corollary}[Sparse mixed-integer quadratic programs]\label{cor:sp:qp}
Let $F_0(x)=\frac12x^{\mathsf T}Hx+b^{\mathsf T}x+c_0$ be rational on a mixed
box $X$ as in \cref{def:sp:instance}, with a tree decomposition of bag size
$p$ containing $\{i,k\}$ whenever $H_{ik}\ne0$, and let $L>0$ be rational
with $L\ge\max_iH_{ii}$ (verified). There is a power of two $M$, computed
from the base instance with $\log_2M\le\poly(I)$, such that for the ambient
perturbation model with marginal law $U_{\sigma,M}$:
\begin{enumerate}[label=(\roman*)]
\item on every draw the algorithm returns a rational global minimizer of
$F_\gamma$ on $X$ and the optimal value, of bit length $\poly(I)$;
\item the expected bit work is at most
$C_0^p[4+(1+n/2)Lw_{\max}/(2\sigma)]^p\poly(I)$.
\end{enumerate}
No growth, uniqueness, convexity, bound on negative inertia, bound on variable
occurrences, or bound on the integer dimension is used.
\end{corollary}

The algorithm is that of \cref{def:sp:algorithm} with two changes. After a
patch output, the strongly convex quadratic program
$\min\{F_\gamma(\bar x_S,y):y\in P\}$ is solved exactly; its unique solution
is rational of polynomial bit length and is computed in polynomial time by
the ellipsoid method for convex quadratic programming
\citep{kozlov1980-the-polynomial-solvability-of-convex}. The fallback is the elementary
enumeration of integer assignments and continuous faces of
\cref{lem:sp:faces}, with $B=3^{n_c}N_Z$; it needs no real algebraic
geometry. The proof is in \cref{app:sp:qp}. For a quadratic $\kappa_3$ may be
replaced by zero in \eqref{eq:sp:matrix}, because the Hessian is constant.

The bound is fixed-width polynomial, as before. It does not contradict the
strong NP-hardness of box-constrained quadratic programming at treewidth two
\citep{pia2026-treewidth-and-the-complexity-of}: the corollary solves the sampled instance,
and its bound is polynomial only when $Lw_{\max}/\sigma$ is polynomially
bounded, a noise level at which the decision thresholds of such hardness
reductions are not preserved (\cref{prop:lim:threshold}).

\subsection{Relation to earlier work}\label{sec:sp:prior}

The ingredients are classical. Dynamic programming over tree decompositions
is nonserial dynamic programming \citep{bertele1972-nonserial-dynamic-programming}. The allowance
$Lw^2/8$ per coordinate is the maximal separation of the
$\alpha$BB underestimator \citep{adjiman1998-a-global-optimization-method-bb}, used here with the
opposite sign as a lower-bound correction. Generic uniqueness and quadratic
growth of minimizers under linear tilts on compact semialgebraic sets are
known qualitatively \citep{lee2017-generic-properties-for-semialgebraic-programs}; they give no finite-law probability
bound. Smoothed analyses of discrete optimization use isolation of winners
\citep{beier2006-typical-properties-of-winners-and,roglin2007-smoothed-analysis-of-integer-programming}; a continuous box has no
positive gap between distinct feasible values, which is why the count above
uses near-optimal grid nodes instead. Treewidth-based linear programming
approximations of sparse polynomial optimization \citep{bienstock2018-lp-formulations-for-polynomial-optimization}
give accuracy-dependent formulations rather than exact optimizers. Exact
quadratic optimization on a box is polynomial on forests and strongly
NP-hard at treewidth two \citep{pia2026-treewidth-and-the-complexity-of}. The new element of
this section is the composition: an expected count of retained bag cells
under one finite independent-noise law that does not grow with the mesh,
combined with sound closure on a certified strongly convex patch, an exact
same-draw fallback, and an output whose evaluation costs polynomial time in
the requested precision.
